Before introducing the hybrid model, we use the SIT cells to characterise the
problem landscape: we establish the DE-LSTM baseline, show that the standard
one-step evaluation protocol saturates on these data, quantify the
cell-to-cell variability that limits pooled prediction, and examine how
training-run variance confounds cross-condition comparisons.

\subsection{Baseline Performance on the SIT Cells}\label{sec:baseline}

Figure~\ref{fig:sit_overview} shows the SOH trajectories for all 20 cells.
Within G1 (identical conditions), final SOH spans 0.067 to 0.885 (a
13.2$\times$ ratio), illustrating the scale of cell-to-cell variability.
G2 cells (elevated temperature) show consistently high residual SOH
(0.660 to 0.924), suggesting the 40\,\textcelsius\ chamber protocol did not
accelerate degradation within the test duration.
G3 cells (elevated temperature and 2C rate) show more diverse outcomes
(0.461 to 1.075), including two cells with moderate to severe degradation.

\begin{figure*}[htbp]
\centering
\includegraphics[width=\textwidth]{fig2_sit_overview}
\caption{SOH trajectories for all 20 SIT cells.}
\label{fig:sit_overview}
\end{figure*}

Table~\ref{tab:sit_baseline} reports DE-LSTM RMSE across all 20 SIT cells
at 30\%, 50\%, and 70\% training splits; values are mean\,$\pm$\,std across the
cells in each group, averaged over two evaluation seeds.

\begin{table}[htbp]
\caption{DE-LSTM baseline RMSE by operating-condition group.}
\label{tab:sit_baseline}
\footnotesize
\begin{tabular*}{\columnwidth}{@{\extracolsep{\fill}}llrrr}
\toprule
Group & Condition & RMSE (30\%) & RMSE (50\%) & RMSE (70\%) \\
\midrule
G1 & 25\,\textcelsius, 1C & $0.030\pm0.025$ & $0.012\pm0.005$ & $0.010\pm0.006$ \\
G2 & 40\,\textcelsius, 1C & $0.018\pm0.021$ & $0.011\pm0.007$ & $0.011\pm0.008$ \\
G3$^\dagger$ & 40\,\textcelsius, 2C & $0.020\pm0.021$ & $0.011\pm0.008$ & $0.010\pm0.008$ \\
\midrule
All$^\dagger$ & {}                   & $0.025\pm0.023$ & $0.011\pm0.006$ & $0.010\pm0.007$ \\
\bottomrule
\multicolumn{5}{@{}l}{\footnotesize $^\dagger$ Excluding the recovery-artefact cell 003-5, whose 30\%-split}\\
\multicolumn{5}{@{}l}{\footnotesize RMSE (0.234) is dominated by its SOH$>$1 segment (50/70\%: 0.017/0.014).}
\end{tabular*}
\end{table}

Two patterns are immediately evident.
First, within G1, per-cell RMSE still varies ${\sim}5\times$ (0.004 to 0.021
at the 50\% split) despite identical conditions
(Figure~\ref{fig:sit_baseline_percell}), and the worst cells at every split
are the catastrophic degraders and the recovery artefact, whose early history
is unrepresentative of their future.
Second, once half a cell's history is available the per-cell errors are small
(${\sim}$0.01) and nearly uniform across operating conditions, a first hint
that the one-step task itself is close to its floor, which the next comparison
makes explicit.

\begin{figure}[htbp]
\centering
\includegraphics[width=\columnwidth]{fig3_sit_baseline}
\caption{DE-LSTM per-cell RMSE at the 50\% split for all 20 SIT cells, sorted by
final SOH.}
\label{fig:sit_baseline_percell}
\end{figure}

\subsubsection{Comparison with Established Methods}

To confirm that this error floor reflects the dataset rather than a weakness of
DE-LSTM specifically, we compare six single-cell predictors on the
17 test cells under the identical protocol: four LSTM-family methods (a
fixed-hyperparameter Simple LSTM, the attention-based
PA-LSTM~\cite{qu2019neural}, GA-LSTM~\cite{ponnambalam2024battery},
and DE-LSTM~\cite{ponnambalam2026delstm}) and two modern non-recurrent
architectures: a temporal convolutional network (TCN, dilated causal
convolutions) and a Transformer encoder, each under a fixed configuration
matched to the Simple-LSTM protocol.
Table~\ref{tab:method_comparison} reports the mean RMSE over the test cells at each split.

\begin{table}[htbp]
\caption{Single-cell predictor and trivial-baseline RMSE on the SIT test cells.}
\label{tab:method_comparison}
\footnotesize
\begin{tabular*}{\columnwidth}{@{\extracolsep{\fill}}lrrr}
\toprule
Method & RMSE (30\%) & RMSE (50\%) & RMSE (70\%) \\
\midrule
\multicolumn{4}{@{}l}{\textit{Recurrent (LSTM family)}} \\
Simple LSTM & $0.0661$ & $0.0352$ & $0.0239$ \\
PA-LSTM~\cite{qu2019neural} & $0.0158$ & $0.0263$ & $0.0173$ \\
GA-LSTM~\cite{ponnambalam2024battery} & $0.0633$ & $0.0198$ & $0.0122$ \\
DE-LSTM~\cite{ponnambalam2026delstm} & $0.0379$ & $0.0118$ & $0.0105$ \\
\midrule
\multicolumn{4}{@{}l}{\textit{Non-recurrent (modern)}} \\
TCN         & $0.1078$ & $0.1017$ & $0.1089$ \\
Transformer & $0.1681$ & $0.1130$ & $0.1001$ \\
\midrule
\multicolumn{4}{@{}l}{\textit{Trivial baselines}} \\
Persistence ($\hat{y}_t=y_{t-1}$) & $0.0123$ & $0.0128$ & $0.0134$ \\
Drift ($\hat{y}_t=2y_{t-1}-y_{t-2}$) & $0.0210$ & $0.0219$ & $0.0228$ \\
Linear AR & $\mathbf{0.0068}$ & $\mathbf{0.0052}$ & $\mathbf{0.0052}$ \\
Cubic curve fit & $1.499$ & $0.158$ & $0.088$ \\
\bottomrule
\end{tabular*}
\end{table}

\begin{figure}[htbp]
\centering
\includegraphics[width=\columnwidth]{fig_method_comparison}
\caption{Mean RMSE of six single-cell predictors and the proposed hybrid across the 30\%, 50\% and 70\% splits.}
\label{fig:method_comparison}
\end{figure}

Table~\ref{tab:method_comparison} shows that the one-step protocol is
saturated on these data: a \textbf{20-cycle linear autoregressor attains the
lowest error at every split}, on par with or ahead of all six deep models, and
even a \textbf{memoryless persistence predictor} ($\hat{y}_t=y_{t-1}$, RMSE
0.013 at 50\%) matches the best deep model (DE-LSTM, 0.0118) and beats the
proposed hybrid (0.0156) on this metric. Among the deep models themselves no
consistent ranking emerges (PA-LSTM leads at
30\%, DE-LSTM at 50\% and 70\%, in sharp contrast to the NASA benchmark where
DE-LSTM reduces RMSE by \mbox{85 to 91\%} over PA-LSTM and
GA-LSTM~\cite{ponnambalam2026delstm}).
This is a property of the task rather than a weakness of any one model, and it
is precisely why, on large-format cells, a predictor must be judged by the
capabilities it provides rather than by its one-step rank, the case the rest of
this paper makes for the proposed hybrid.
The explanation is that rolling one-step-ahead prediction from \emph{observed}
windows is a near-persistence task on smooth SOH trajectories: successive
cycles differ by a fraction of a percent, the predictor always receives the
true recent history, and a linear map over the last 20 observations is
essentially optimal.
The task never requires anticipating a knee or a crash, only following one
cycle behind it, which is also why the curve-fitting baseline, the one
trivial method forced to extrapolate, fails by an order of magnitude.
We term this \textbf{benchmark saturation}: the protocol under which this
method lineage, and much of the field, ranks models cannot in fact
distinguish them.
A two-way variance decomposition over the six deep predictors makes the point
quantitatively: the identity of the \emph{cell} explains 58\%/31\%/25\% of
prediction-error variance at the 30/50/70\% splits respectively, whereas the
choice of \emph{method} explains only 15\% to 17\% at every split, with the
remainder attributable to cell-method interaction and training-run variance.
Model selection should therefore rest not on one-step leaderboards but on
capabilities the task actually demands: predicting cells with \emph{no}
usable history, quantifying uncertainty honestly, and transferring across
operating regimes and chemistries, the capabilities evaluated in
Section~\ref{sec:results}.

\subsection{Cell-to-Cell Variability and the Two-Sided Failure Mode}\label{sec:variability}

We examine this cell-to-cell variability from three angles: how predictable each
cell is on its own, what happens when a model is pooled across cells, and whether
early-life data can identify the problem cells in advance.

\subsubsection{Variability at Matched Exposure}
The 13.2$\times$ spread in final SOH (0.067 to 0.885) compares cells at different
points in their lives, because they reach end of life after different cycle counts.
To measure the variability fairly we compare the G1 cells at matched exposure. At a
common cycle count, the shortest G1 record of 684 cycles, SOH ranges from 0.20 to
0.90, a 7.6$\times$ ratio; at a common cumulative discharge throughput of
26.8\,kAh it ranges from 0.07 to 0.92, a 13.7$\times$ ratio. The spread therefore
remains roughly 8 to 14 times under every matched-exposure definition, so it is not
an artefact of unequal test durations but a genuine dispersion among
nominally identical cells.

\subsubsection{Within-Cell Predictability}

A cubic polynomial fit to each G1 cell's SOH trajectory achieves a mean
$R^2$ of 0.964 (minimum: 0.908 for cell 001-6), confirming that individual
trajectories are highly predictable given sufficient training data.
On average 3.6\% of the trajectory variation is not captured by this
particular cubic fit; this residual reflects the choice of fit and is not
an irreducible noise floor.
The challenge is therefore not within-cell prediction but
\textit{between-cell diversity}: knowing a cell's past trajectory does not
reveal where it sits in the manufacturing distribution.

\subsubsection{Cross-Cell LOOCV: A Two-Sided Failure Mode}

\begin{figure*}[htbp]
\centering
\includegraphics[width=\textwidth]{fig4_loocv}
\caption{G1 LOOCV: same-cell vs cross-cell RMSE at the 50\% split.}
\label{fig:loocv}
\end{figure*}

Leave-one-out cross-validation (LOOCV) on the 10 G1 cells reveals a nuanced
picture (Table~\ref{tab:loocv}, whose factor column gives the
cross-cell/same-cell RMSE ratio, with values below 1 favouring cross-cell).
The pooled cross-cell model (a plain LSTM, $H=64$, $W=20$, trained on the
\emph{full} trajectories of
the other nine cells and tested on the held-out cell) achieves a lower
\textit{mean} RMSE than same-cell training (which uses only the first 50\% of the
held-out cell; 0.049 vs.\ 0.060 at 50\% split), but this average conceals
opposing failure modes:

\begin{table}[htbp]
\caption{Same-cell vs cross-cell RMSE for the 10 G1 cells (50\% split).}
\label{tab:loocv}
\begin{tabular*}{\columnwidth}{@{\extracolsep{\fill}}lrrrr}
\toprule
Cell & SOH$_\text{end}$ & Same-cell & Cross-cell & Factor \\
\midrule
001-1 & 0.767 & 0.0363 & 0.0197 & 0.54$\times$ \\
001-2 & 0.878 & 0.0149 & 0.0262 & 1.77$\times$ \\
001-3 & 0.493 & 0.0938 & 0.0628 & 0.67$\times$ \\
001-4 & 0.709 & 0.0348 & 0.0389 & 1.12$\times$ \\
001-5 & 0.885 & 0.0166 & 0.0340 & 2.04$\times$ \\
001-6 & 0.854 & 0.0153 & 0.0189 & 1.24$\times$ \\
001-7 & 0.885 & 0.0162 & 0.0311 & 1.92$\times$ \\
001-8 & 0.067 & 0.1746 & 0.1183 & 0.68$\times$ \\
101-1 & 0.696 & 0.0424 & 0.0425 & 1.00$\times$ \\
101-3 & 0.118 & 0.1553 & 0.1015 & 0.65$\times$ \\
\midrule
Mean  &       & 0.0600 & 0.0494 & 0.82$\times$ \\
\bottomrule
\end{tabular*}
\end{table}

The cross-cell model is better for 4 of 10 cells, and the largest gains are
on the catastrophic degraders (001-8: same-cell RMSE = 0.175, cross-cell =
0.118; 101-3: 0.155 vs.\ 0.102), where same-cell training on the first 50\%
of a catastrophically failing cell provides a severely biased trajectory
estimate.
But pooling degrades five of the remaining cells by up to 2$\times$ (001-5:
0.017 same-cell vs.\ 0.034 cross-cell; 001-2 and 001-7 are also harmed by
factors of 1.8 to 1.9), and nothing in a cell's early data reveals in advance
whether it will be helped or harmed: the harmed cells are the \emph{most
robust} cells in the population, whose own smooth history is a better guide
than the pooled average.
The split is not sensitive to the training budget: when the cross-cell model
is retrained on only the first 50\% of each source cell (matching the
same-cell information budget), the mean falls to 0.019 vs.\ 0.023 same-cell
and the helped/harmed split remains 6/10, so the two-sided pattern is a
property of pooling itself, not of the data advantage.

Figure~\ref{fig:loocv} makes this pattern visual: the left panel contrasts
same-cell and cross-cell RMSE for each G1 cell, while the right panel plots each
cell's degradation factor against its final SOH, showing cross-cell training
helping the low-SOH catastrophic degraders but hurting the high-SOH robust cells.
This \textbf{two-sided failure mode}, where cross-cell training rescues
catastrophic degraders but unpredictably harms a minority of ordinary cells,
motivates cell-specific adaptation via the fingerprint encoder in
Section~\ref{sec:solution}.

\subsubsection{Early-Life Fingerprint}

The practical value of this insight depends on whether the early SOH
trajectory (observable in the first weeks of deployment) can identify
which cells are likely catastrophic degraders.
Table~\ref{tab:fingerprint} reports Pearson correlations between early-life
statistics and final SOH across 47 MIT LFP cells (SOH$_\text{end}<1$,
excluding 2 capacity-recovery artefacts).

\begin{table}[htbp]
\caption{Correlation of early-life features with final SOH.}
\label{tab:fingerprint}
\begin{tabular*}{\columnwidth}{@{\extracolsep{\fill}}llrrl}
\toprule
$K$ cycles & Feature & $r$ & $p$ & Sig. \\
\midrule
10 & early\_std & $+$0.322 & 0.027 & * \\
10 & max\_drop  & $+$0.371 & 0.010 & * \\
10 & slope      & $-$0.073 & 0.625 &   \\
30 & slope      & $-$0.297 & 0.043 & * \\
50 & slope      & $-$0.380 & 0.008 & ** \\
\bottomrule
\multicolumn{5}{@{}l}{\footnotesize Significance: $^{**}\,p<0.01$; $^{*}\,p<0.05$.}
\end{tabular*}
\end{table}

The 50-cycle slope ($r=-0.38$, $p=0.008$) and 10-cycle max\_drop ($r=+0.37$,
$p=0.010$) are the most predictive individual features
(Figure~\ref{fig:fingerprint}), but with $R^2 \approx 0.14$
these statistics explain only 14\% of the variance in final SOH.
Linear early-life statistics are insufficient for reliable individual-cell
prediction; this motivates the CNN fingerprint encoder, which learns a
nonlinear representation from the full 30-cycle trajectory rather than
hand-crafted summary statistics.

\begin{figure*}[htbp]
\centering
\includegraphics[width=\textwidth]{fig10_fingerprint}
\caption{Early-life predictors of final SOH: 50-cycle slope and 10-cycle maximum capacity drop.}
\label{fig:fingerprint}
\end{figure*}

\subsection{Cross-Condition Transfer Is Dominated by Training-Run
Variance}\label{sec:brittleness}

\begin{table}[htbp]
\caption{Cross-condition transfer RMSE between operating-condition groups.}
\label{tab:brittleness}
\resizebox{\columnwidth}{!}{%
\begin{tabular}{llrr}
\toprule
Train & Test & RMSE & Factor vs G1 LOOCV \\
\midrule
G1 (LOOCV) & G1          & $0.062 \pm 0.065$ & $1.0\times$ (baseline) \\
G1          & G2          & $0.008 \pm 0.002$ & $0.1\times$ \\
G1          & G3          & $0.034 \pm 0.041$ & $0.6\times$ \\
G2          & G3          & $0.145 \pm 0.139$ & $2.4\times$ \\
G1+G2       & G3          & $0.145 \pm 0.133$ & $2.4\times$ \\
\bottomrule
\end{tabular}%
}
\end{table}

Table~\ref{tab:brittleness} reports the cross-condition transfer RMSE
(mean\,$\pm$\,std across the target cells), with the degradation factor giving
each transfer's RMSE relative to the within-distribution G1 LOOCV baseline;
Figure~\ref{fig:brittleness} plots the same values, with the dashed line marking
that baseline.
Two confounds emerge, and together they carry a methodological warning.
First, G1$\rightarrow$G2 (temperature shift only) yields a far lower RMSE
(0.008) than the G1 LOOCV baseline (0.062), not because temperature
transfer is easy, but because G2 cells happen to have smooth, high SOH
trajectories well within the G1 training distribution: cross-condition RMSE
comparisons are confounded by the difficulty distribution of the test set.
Second, and more consequentially, the transfer factors are dominated by
\emph{training-run variance} rather than by the condition shift itself.
Each transfer row is one pooled model, and independent training draws of the
same experiment produce qualitatively different conclusions: an earlier run
of G1$\rightarrow$G3 with library-default (unseeded) initialisation yielded
$2.3\times$ degradation, whereas the deterministic pipeline reported here
yields $0.6\times$, i.e.\ \emph{better} than within-distribution.
Even within Table~\ref{tab:brittleness} the instability is visible:
training on G1 alone transfers to G3 at 0.034, yet \emph{adding} G2 data
(G1+G2$\rightarrow$G3) worsens the same transfer to 0.145, a difference no
condition-shift mechanism explains.
At one-step granularity, cross-condition effects on these data are smaller
than the noise floor of neural-network training itself, another face of the
benchmark saturation established in Section~\ref{sec:baseline}.
We therefore do not claim a cross-condition brittleness effect; the
defensible claims are the test-set difficulty confound and the dominance of
run variance, both of which argue for capability-based evaluation with
uncertainty quantification rather than single-model RMSE comparisons.

The within-distribution baseline used here (G1 LOOCV, $0.062$) trains on the
first 50\% of each source-group cell, identical to the cross-condition
transfers, so every degradation factor in Table~\ref{tab:brittleness} is
computed within one consistent protocol.
It is therefore larger than the full-trajectory cross-cell LOOCV reported in
Section~\ref{sec:variability} ($0.049$), which trains on the complete
trajectories of the other nine cells and so sees their degradation tails; the
two baselines answer different questions and should not be compared directly.

\begin{figure}[htbp]
\centering
\includegraphics[width=\columnwidth]{fig5_brittleness}
\caption{Mean RMSE for each cross-condition transfer under the deterministic pipeline.}
\label{fig:brittleness}
\end{figure}

\subsection{Requirements for a Robust Predictor}\label{sec:motivation}

The preceding results expose why no single-cell predictor stands out on SIT and
what a better model must instead provide.
The limiting factor is not the sequence model or its hyperparameters but
cell-to-cell variability.
Under nominally identical G1 conditions, per-cell optimised DE-LSTM RMSE still
varies ${\sim}5\times$ across the ten cells
(Section~\ref{sec:baseline}), because a model trained on the early portion of a
single cell's trajectory cannot tell where that cell sits in the manufacturing
distribution: the first half of a catastrophic degrader closely resembles that
of a moderate one.
No amount of per-cell hyperparameter search resolves this, because the missing
information is not present in the individual cell's early history.

Pooling data across cells, the obvious alternative, introduces a two-sided
failure mode (Section~\ref{sec:variability}).
A single model trained on the whole population helps catastrophic degraders,
for which the population average is a better prior than their own ambiguous
early trajectory, but it degrades half of the ordinary cells by up to
2$\times$, and which cells suffer is not predictable a priori.
Neither per-cell tuning nor naive pooling addresses the underlying problem.

These observations, together with the saturation of the one-step metric,
reframe the task.
What matters is not squeezing the last fraction of a percent from a saturated
leaderboard but providing capabilities the leaderboard does not measure: a
degradation prior broad enough to recognise extreme trajectories, prediction
for cells with little or no usable history, and honest uncertainty estimates.
The hybrid model provides the first through pretraining on a large, diverse LFP
corpus (Section~\ref{sec:datasets}), the second through an early-life
fingerprint encoder that enables prediction without target-cell fine-tuning,
and the third through MC-Dropout with conformal calibration.
